\documentclass[11pt]{article}

\usepackage[a4paper,margin=2.7cm]{geometry}
\usepackage{amsmath,amssymb,amsthm,mathtools}
\usepackage{hyperref}

\newtheorem{theorem}{Theorem}[section]
\newtheorem{proposition}[theorem]{Proposition}
\newtheorem{lemma}[theorem]{Lemma}
\newtheorem{remark}[theorem]{Remark}

\newcommand{\C}{\mathbb C}
\newcommand{\R}{\mathbb R}
\newcommand{\dd}{\,d}
\newcommand{\1}{\mathbf 1}

\title{The DtN Zero Condition as a Boundary-Curve Equation}
\author{}
\date{}

\begin{document}
\maketitle

\section{The affine DtN map}

Let \(U\subset\C\simeq\R^2\) be a bounded domain with real-analytic boundary, and
let \(\rho:[0,\infty)\to\R\) be real analytic.  For boundary data
\(\varphi\), let \(w_\varphi\) solve
\begin{equation}\label{eq:dirichlet-problem}
  -\Delta w_\varphi=\rho(|z|)
  \qquad\text{in }U,\qquad
  w_\varphi|_{\partial U}=\varphi .
\end{equation}
The corresponding inhomogeneous Dirichlet-to-Neumann map is the affine map
\[
  \mathcal N_U^\rho(\varphi):=\partial_\nu w_\varphi|_{\partial U},
\]
where \(\nu\) is the outward unit normal.

Thus
\[
  \mathcal N_U^\rho(0)=0
\]
means exactly that the zero-Dirichlet solution \(u:=w_0\) satisfies the
overdetermined boundary conditions
\begin{equation}\label{eq:overdetermined}
  -\Delta u=\rho(|z|)\quad\text{in }U,\qquad
  u=0,\quad \partial_\nu u=0\quad\text{on }\partial U .
\end{equation}
Equivalently,
\[
  \nabla u=0\qquad\text{on }\partial U,
\]
because the tangential derivative of the zero trace also vanishes.

\begin{remark}[Flux compatibility]
For the nonsingular problem \eqref{eq:dirichlet-problem}, the condition
\(\mathcal N_U^\rho(0)=0\) implies
\[
  \int_U \rho(|z|)\dd A=0
\]
by integrating the equation.  In particular, if \(\rho>0\) in \(U\), this
condition is impossible.  In localization applications one often has either a
sign-changing \(\rho\), or an additional radial point-mass term at the origin.
The boundary-curve calculation below is unchanged by a radial point mass,
because the angular derivative annihilates it.
\end{remark}

\section{Boundary jets}

Let \(\gamma(s)\) be an arclength parametrization of a boundary component
\(\Gamma\subset\partial U\), and let
\[
  T(s):=\gamma'(s)
\]
be the unit tangent.  Let \(\nu(s)\) be the outward unit normal.  The condition
\(\nabla u=0\) on \(\Gamma\) implies the following second-jet identities.

\begin{lemma}\label{lem:hessian}
On \(\Gamma\),
\[
  D^2u(T,T)=0,\qquad D^2u(T,\nu)=0,\qquad D^2u(\nu,\nu)=-\rho(|\gamma|).
\]
Consequently,
\begin{equation}\label{eq:hessian-normal}
  D^2u=-\rho(|\gamma|)\,\nu\otimes\nu
  \qquad\text{on }\Gamma .
\end{equation}
\end{lemma}

\begin{proof}
Since \(u(\gamma(s))=0\), differentiating once gives
\[
  \nabla u(\gamma(s))\cdot T(s)=0.
\]
This is already contained in \(\nabla u=0\).  Differentiating again,
\[
  D^2u(T,T)+\nabla u\cdot T'=0.
\]
The second term vanishes on \(\Gamma\), hence \(D^2u(T,T)=0\).

Similarly, differentiating \(\partial_\nu u=\nabla u\cdot\nu=0\) tangentially gives
\[
  D^2u(T,\nu)+\nabla u\cdot\nu'=0,
\]
so \(D^2u(T,\nu)=0\).  Finally,
\[
  -\Delta u
  =
  -D^2u(T,T)-D^2u(\nu,\nu)
  =
  \rho(|\gamma|),
\]
and therefore \(D^2u(\nu,\nu)=-\rho(|\gamma|)\).
\end{proof}

\section{Angular derivative and the boundary ODE}

Let
\[
  L:=x\partial_y-y\partial_x
\]
be the angular derivative, and set
\[
  v:=Lu.
\]
Since \(L\) commutes with \(\Delta\) and annihilates every radial function,
\[
  -\Delta v=L(-\Delta u)=L(\rho(|z|))=0
  \qquad\text{in }U .
\]
Moreover \(v=0\) on \(\partial U\), because \(\nabla u=0\) there.  Hence, by
uniqueness for the Dirichlet problem,
\[
  v\equiv0\qquad\text{in }U,
\]
and therefore
\[
  \partial_\nu v=0\qquad\text{on }\partial U .
\]

Now compute \(\partial_\nu v\) directly on \(\Gamma\).  Write
\[
  X(z):=(-y,x)
\]
so that \(Lu=X\cdot\nabla u\).  Since \(\nabla u=0\) on \(\Gamma\),
\[
  \partial_\nu v
  =
  D^2u(\nu,X)
  =
  -\rho(|\gamma|)(\nu\cdot X(\gamma))
\]
by \eqref{eq:hessian-normal}.  Since \(\nu\cdot X(\gamma)=\pm\gamma\cdot T\),
depending only on the orientation convention for \(T\), we obtain the intrinsic
boundary equation
\begin{equation}\label{eq:boundary-ode}
  \rho(|\gamma(s)|)\,\gamma(s)\cdot\gamma'(s)=0 .
\end{equation}
Equivalently,
\begin{equation}\label{eq:radius-ode}
  \rho(|\gamma(s)|)\frac{\dd}{\dd s}|\gamma(s)|^2=0 .
\end{equation}

\begin{theorem}\label{thm:dtn-to-circle}
Assume \(\rho\) is real analytic and not identically zero.  If
\(\mathcal N_U^\rho(0)=0\), then every boundary component of \(U\) is a circle
centered at the origin, except possibly components lying entirely on a zero
circle of \(\rho\); those are also centered circles.  In particular, if
\(U\) is bounded and simply connected, then \(U\) is a disk centered at the
origin.
\end{theorem}

\begin{proof}
Equation~\eqref{eq:radius-ode} says that the product of the two real-analytic
functions
\[
  s\mapsto \rho(|\gamma(s)|),
  \qquad
  s\mapsto \frac{\dd}{\dd s}|\gamma(s)|^2
\]
vanishes identically on the connected parameter circle.

If \(\rho(|\gamma(s)|)\) is not identically zero, then its nonzero set is open
and dense.  On that set, \(\frac{\dd}{\dd s}|\gamma(s)|^2=0\), and by analyticity
this derivative vanishes for all \(s\).  Hence \(|\gamma|\) is constant.

If instead \(\rho(|\gamma(s)|)\equiv0\), then the connected set
\(\{|\gamma(s)|:s\in\R/2\pi\mathbb Z\}\) is contained in the zero set of the
real-analytic function \(\rho\).  Since \(\rho\not\equiv0\), its zero set has no
intervals.  Therefore \(|\gamma|\) is again constant.

Thus each boundary component is a level set \(|z|=R\), hence a circle centered
at the origin.  A bounded simply connected domain with one centered circular
boundary component is the disk bounded by that circle.
\end{proof}

\begin{remark}[Radial singularities]
If the interior equation is instead
\[
  -\Delta u=\rho(|z|)-c\,\delta_0,
\]
with \(0\in U\), the same argument applies on \(U\setminus\{0\}\) and
distributionally across \(0\), because \(L\delta_0=0\).  Thus \(v=Lu\) is still
harmonic with zero boundary values, and the same boundary equation
\eqref{eq:boundary-ode} follows.
\end{remark}

\end{document}
